\section{Introduction}
Reward functions translate a control objective into the signal optimized by a reinforcement-learning (RL) policy. A reward that appears reasonable can still suppress exploration, overemphasize a proxy, or create a policy that exploits the reward without completing the intended task~\cite{sutton2018rl,ng1999policy,hadfield2017inverse,amodei2016concrete}. Human reward engineering therefore proceeds as a loop: design a reward, train a policy, inspect the resulting behavior, and revise the objective. The policy-training step is expensive, while the information produced by a failed run is often richer than its final score.

Code-capable LLMs have made executable reward synthesis possible~\cite{language2rewards,text2reward,eureka}. Most automated systems explore multiple reward candidates, refine candidates with feedback, or organize reward search as an agentic process~\cite{card,rfagent,revolve}. These advances establish that an LLM can participate in reward design. Our focus is narrower and operational: after one generated reward produces a poor policy, can an autonomous agent use that exact training run to continue improving the same reward program?

We propose CREATE, a \emph{Closed-loop Reward Evolution Agent with Training--Evaluation Feedback}. CREATE maintains one active reward lineage. At each iteration it trains a policy with the generated reward, evaluates the policy with the environment-native task objective, summarizes reward-component and behavioral evidence, and chooses a bounded revision. A lineage memory records what changed and what happened next; a best archive prevents a later unsuccessful revision from erasing an earlier solution. This supervisory reward-design loop complements the lower-level policy learner and directly fits learning-based control and autonomous decision-making.

The paper makes two contributions. First, it formulates LLM reward design as a persistent training-feedback loop with explicit separation between the policy-training reward and the native evaluation objective, plus lineage memory and best-solution retention. Second, it reports multi-seed evidence on discrete- and continuous-action control tasks, including complete seed-wise evolution, representative revision traces, and an unconstrained-revision ablation.

\section{Related Work}
Classical reward shaping studies how auxiliary objectives affect policy learning; potential-based shaping provides invariance guarantees under specific transformations~\cite{ng1999policy}, while reward machines expose task structure explicitly~\cite{rewardmachines}. These methods remain dependent on a designer-specified representation.

Recent work uses LLMs to generate reward code from task descriptions and environment interfaces. Language to Rewards and Text2Reward translate language into executable rewards~\cite{language2rewards,text2reward}; Eureka couples code generation with evolutionary search and reward reflection~\cite{eureka}. CARD introduces dynamic feedback, RF-Agent uses language-agent tree search, and REvolve incorporates human feedback~\cite{card,rfagent,revolve}. General self-refinement and reflection methods motivate persistent language-agent state~\cite{selfrefine,reflexion}. Most closely related, diagnostic-driven refinement studies how training failures can guide reward repair and identifies limitations in continuous control~\cite{wang2026diagnostic}. CREATE does not claim that a single lineage universally dominates population search. It studies a compact autonomous loop in which evidence from one reward-policy pair is retained and used to revise that same reward program.
